\documentstyle[12pt,newsymb]{article}%[11pt]{block}
%\def\thesection{\Roman{section}}
%\textwidth=6in
%\advance\hoffset -2.5truein
%\renewcommand{\theequation}{\thesection.\arabic{equation}} %set counter at
%beginning of each section`\setcounter{equation}{0}'
%\textheight %text height in inches
%\topmargin % 0 is the default which gives a one inch margin
%\headheight % about 1/4 in is the default, 0 removes
%\headsep % the distance between the head and the body of the text the
%default is about 1/4 in
%\pagestyle{normal} %{empty}
\def\negate#1{\mbox{$\;$\begin{picture}(10,12)
\put(0,-2){\line(1,1){10}}
\put(6,0){\makebox(0,0)[b]{#1}}
\end{picture}$\;$}}
%\advance\hoffset by .5truein\relax%controls horizontal position on the page
\input bigc
\newtheorem{example}{Example}[section]
\newtheorem{corollary}{Corollary}[section]
\newtheorem{lemma}{Lemma}[section]
\newtheorem{theorem}{Theorem}[section]
\newtheorem{proposition}{Proposition}[section]
\newtheorem{definition}{Definition}[section]
\newtheorem{notations}{Notations}[section]
\newtheorem{assumption}{Assumption}[section]
\newtheorem{remark}{Remark}
\newtheorem{problem}{Problem}[section]
\begin{document}
%\baselineskip=18pt

\title{Polynomials with non-negative coefficients}

\author{R. W. Barnard \and W. Dayawansa \thanks{Formerly at Texas Tech
 University,\newline Currently in the Department of Electrical Engineering, University of Maryland,College Park, Md 20742.} \and K. Pearce \and D.Weinberg}
\date{}
\maketitle
\begin{center}{Department of Mathematics, Texas Tech University,
Lubbock, TX 79409}\end{center}

\section{Introduction}
{\em Can a conjugate pair of zeros be factored from a polynomial
with nonnegative coefficients so that the resulting polynomial still
has nonnegative coefficients?}

This question has attracted considerable attention during the last
few years because of its seemingly elementary nature and its
potential for applications in number theory and control theory.  A
proposed answer to this question arose as a conjecture out of the
work in \cite{5}, where explicit bounds were determined for the
constants occurring in the Beauzamy-Enflo generalization \cite{1},
\cite{2} of Jensen's Inequality.  An affirmative answer to this
question was conjectured, independently, by B. Conrey in connection
with some of his work in number theory.  Conrey announced the
conjecture at the annual West Coast Number Theory Conference held in
December 1987.

The main theorem in this note gives a positive answer to the
question.  This will be proved in section 2.  The principle
ingredients of the proof are an idea from index theory, classical
properties of polynomials and a significant lemma (Lemma 2.1) which
we prove in section 3.  This lemma states some strong consequences
for the case of equality holding in Descartes' Rule of Signs (see
\cite{4}).  We also prove a corollary of the main theorem which
describes the region into which certain zeros can be moved while
preserving the nonnegativity of the coefficients.


Consider the polynomial defined by $P_N(z) = 1 + z^N$.  We note that


\begin{equation} P_N(z) = \prod_{l=0}^{N-1}\left(1 -
\frac{z}{e^{i\left(\frac{\pi}{N} +\frac{2\pi l}{N}\right)}}\right) =
(1 - 2\cos\theta_0z +z^2)\sum_{k=0}^{N-2} b_kz^k\label{eq:1}\end{equation} has
$b_{k} >0, 0\leq k\leq N-2$, if $\theta_0 = \pi/N$, while any
other choice for $\theta_0$ produces a factor with some negative
$b_k$ coefficients.

Thus, one initial suggestion for the general question of factoring
out a conjugate pair of zeros was to factor out a pair of zeros with
greatest real part.  The authors have used fairly straightforward
arguments to show that if the degree of the polynomial is less than
or equal to 5, then a conjugate pair of zeros of greatest real part
can be factored out and the resulting polynomial will still have
nonnegative coefficients.  However, if $$p(z) = 140 + 20z + z^2 +
1000z^3 + 950z^4 + 5z^5 + 20z^6$$ with $z_0, \overline{z}_0$
approximately equal to $0.392 + 6.390i, 0.392 - 6.392i$, resp., as
the conjugate pair of zeros of greatest real part, then the resulting
factors are $(z - z_0)(z - \overline{z}_0)$ and $(140.1 + 22.3z -
1.53z^2 + 1000z^3 + 966.7z^4)$.

The other choice, for which zeros to factor out, suggested by the
example $P_N$ in (1), is the pair determined by the zero in the
upper half-plane with smallest positive argument.  Indeed, in this
paper we prove the following:\vspace{12pt}

\begin{theorem} Let $p$ be a polynomial of degree $N$, $p(0) =1 $,
with nonnegative coefficients and with zeros $z_1, z_2,\dots, z_N$.
For $t\geq 0$ write $$p_{t} (z) = \prod_{1\leq j \leq N \atop{ |Arg
z_{j} | > t}} \left( 1 - \frac{z}{z_{j}} \right)$$ Then, if $p_t\neq
p$, all of the coefficients of $p_t$ are positive.
\end{theorem}

\begin{remark} We note that a slightly sharper version
of the Theorem 1.1 is true.  Suppose that the given polynomial $p$
has nonnegative coefficients and that we divide it by a single quadratic
factor corresponding to any conjugate pair of zeros of smallest
argument in magnitude.  Then the quotient polynomial has nonnegative
coefficients.  Furthermore, if the conjugate pair of zeros is unique
 and simple or if $p$ has strictly positive coefficients
then the quotient polynomial will have strictly positive
coefficients.  The example $p(z) = (1 + z^{2})^{2}$ shows the
necessity of these conditions.  The sharpened version of the theorem
under the uniqueness hypothesis follows from the main theorem.  Now
let us consider the case when several complex conjugate pairs of
zeros of $p$, say $\{ z_{i} , \bar{z}_{i} \}_{i=1, \ldots, k},$ have
the smallest argument in magnitude and assume that all
coefficients of $p$ are strictly positive.  A local
perturbation can be made of $z_{1}$ and $\bar{z}_{1}$ to $w_{1}$ and
$\bar{w}_{1}$ such that the argument of $w_{1}$ is less than that of
$z_1 $.  Under the local perturbation the polynomial $p$ changes to
$q$ which also has strictly positive coefficients. By Theorem 1.1, $q
/{(z-w_{1})(z-\bar{w}_{1})} $ has strictly positive coefficients.
Since $p/{(z-z_{1})(z-\bar{z}_{1})} = q/{(z-w_{1})(z-\bar{w}_{1})}$
the desired conclusion follows.  \end{remark}

\begin{corollary} Let $$p(z) = \sum^{N}_{k=0} a_{k} (z_{0}) z^{k} =
(z-z_{0})(z-\bar{z}_{0}) \sum^{N-2}_{k=0} b_{k} z^{k}$$ be a real
polynomial with $a_{k} (z_{0}) \geq 0.$ Let $z_{0}$ be a zero of $p$
in the upper half plane with smallest argument.  If $z_{1}$ is any
complex number such that $$|z_{1}| \geq |z_{0}| \mbox{ and } Re \{
z_{1} \} \leq Re \{ z_{0} \} \ \ \ \ \ \ \ \ \ \ \ $$ then
$$\sum^{N}_{k=0} a_{k} (z_{1} ) z^{k} = (z-z_{1}) (z-\bar{z}_{1})
\sum^{N-2}_{k=0} b_{k} z^{k} $$ has $a_{k} (z_{1}) \geq a_{k}
(z_{0})$. \end{corollary}

\noindent {\bf Proof:} The proof follows directly by noting, for $ z_0  =  r e^{i\theta}=x + iy $,

\begin{eqnarray*}
 \sum^N_{k=0}
a_k(z_0)z^k & = & (r^2-2{r}z\cos{\theta} +z^2)\sum^N_{k=0} b_kz^k \\ &  = & (x^2+y^2 -2xz+z^2)\sum^N_{k=0} b_kz^k\end{eqnarray*}
and comparing coefficients.


It has been recently shown by R. Evans and P. Montgomery\cite{6} that in
the special case when $p$ is the polynomial $P_N$ defined in (1) that
the corresponding reduced $p_t$ polynomials are all strictly unimodal
(in their coefficients).  Also, R. Evans and J. Greene \cite{3} have
shown for polynomials $p$ of the form $p(z)=(z^{Mk}-1)/(z^k-1)$ that
the corresponding reduced $p_t$ polynomials all have positive coefficients.

\section{Proof of theorem 1.1}

We note to prove Theorem 1.1 it suffices, by a successive reduction
argument, to show that if we remove from $p$ a conjugate pair of
zeros of smallest argument in the absolute value, then the resulting
polynomial still has nonnegative coefficients.  


\noindent {\bf Proof.}  

Let

\begin{equation} p(z) = \sum^N_{k=0} a_k z^k = (1 -\frac{2\cos
\theta_0 z}{r_0} + \frac{z^2}{r{_0^2}})\sum^{N-2}_{k=0} b_k
z^k\label{eq:2}\end{equation} 

\noindent with $a_k \geq 0, 0\leq k\leq N$.  We may assume $a_N > 0$.
Let $\{r_\ell e^{i\theta_\ell}\}$ be the zero set of $p$ with $0 <
\theta_0 \leq |\theta_\ell|$ for each $\ell$. We may assume $ r_0 =1$
 in (\ref{eq:2}) by using  $ p(r_{0}z)$. Since the coefficients of $p$ depend
continuously on the zeros of $p$, we may assume that $p$ has only one
zero, say $z_0 = e^{i \theta_0}$, on the ray $\{r e^{i\theta_0}: 0
< r < \infty\}$ and that this zero is simple.  Otherwise, a
sequence of polynomials with this property can be chosen to converge
to the desired polynomial with the appropriate inequalities holding.
 We will prove that $b_k > 0$ for $ 0 \leq k
 \leq {N-2}$. 

Now 
\begin{eqnarray*} \sum^{N-2}_{k=0} b_k z^k &=& \frac{p(z)}{1 -
2\cos \theta_0 z+ z^2}\\ &=& \sum^{\infty}_{k=0} \frac{\sin
(k+1)\theta_0}{\sin \theta_0} z^k \sum^N_{k=0} a_k z^k =
\sum^{N-2}_{k=0} \left(\sum^k_{\ell=0} \frac{\sin (\ell +1)
\theta_0}{\sin \theta_0} a_{k -\ell}\right) z^k.\end{eqnarray*} 

\noindent This gives \begin{equation} b_k = \sum^k_{\ell =0}
\frac{\sin (\ell +1)\theta_0}{\sin \theta_0} a_{k -
\ell}.\label{eq:3}\end{equation} 

\noindent Noting that $$z^N p(1/z) = \sum^N_{k=0} a_{N-k} z^k =
\left(z^2 -2\cos \theta_0 z+1\right) \sum^{N-2}_{k=0} b_{N-2-k} z^k$$ 

\noindent it will suffice to show that $b_k >  0$ for $0\leq k \leq
[\frac{N-2}{2}]$.  From (\ref{eq:3}) if \linebreak $\sin (\ell +1)
\theta_0 > 0$ for $\ell = 0, 1, \cdots, k$, i.e., when $0 \leq
\theta_0 < \pi/(k+1)$ it 

\begin{figure}\vspace{3in}\caption{}\end{figure} 


follows that $b_k > 0$.  Also, if all
the zeros of $p$ are in the closure of the left half plane then the
result is clear since $p$ can be put into the form: $p(z) = \prod
(1 - \frac{2\cos \theta_\ell}{r_\ell} z + \frac{z^2}{r^2_{\ell}}) \prod(1 + \frac{z}{r_m})$.  Thus,
we may assume
\begin{equation} \pi/\left(\left[\frac{N-2}{2}\right] +1\right) =
\pi/\left(\left[\frac{N}{2}\right]\right) \leq \theta_0 <
\pi/2.\label{eq:4}\end{equation}

We will assume that there is a $k$ with

\begin{equation} b_k \leq 0\label{eq:5}\end{equation} 

\noindent to reach a contradiction.  Consider the function defined by
$$F(z) = {z}^{-k-1} p(z) = \sum^N_{\ell =0} a_\ell z^{\ell- k-1}$$
having the same zeros as $p$.

We observe from the hypothesis that $F$ has no zeros between the rays
defined by $\{te^{i\theta_0}: 0 < t < \infty\}$ and $\{te^{-i
\theta_0}: 0 < t < \infty\}$.  Motivated by this, let $\Gamma =
\Gamma ^+\bigcup\overline\Gamma ^+$ be the curve shown in Figure 1,
which is symmetric about the reals with $\Gamma^+ =
\bigcup^{5}_{\ell=1} \Gamma_\ell$ where $\Gamma_1 = \{Re^{i\theta}:0
\leq \theta \leq \theta_0\},\ \Gamma_2 = \{te^{i\theta_0}: 1 + s \leq
t \leq R\},\ \Gamma_3 = \{z_0 + se^{i\theta}: \theta_0 - \pi \leq
\theta \leq \theta_0\},\ \Gamma_4 = \{t e^{i\theta_0}: r \leq t \leq
1 -s\}$ and $\Gamma_5 = \{re^{i\theta}: 0 \leq \theta \leq
\theta_0\}$ and $r, s$, and $1/R$ will be chosen sufficiently small.
We will show that under the assumption (\ref{eq:5}) the index of $F$
with respect to $\Gamma$ about any interior point is positive,
implying the existence of a zero of $F$ inside $\Gamma$.  This would
contradict that $\theta_0$ is the smallest positive argument of the
zeros of $p$ in the upper half plane.

The function $F$ has the form \begin{equation} F(z) = a_0 z^{-k-1} +
\cdots + a_{k+1} + a_{k+2} z + \cdots a_N
z^{N-k-1}.\label{eq:6}\end{equation} From symmetry we can examine
$\Delta \arg F(z)$, the change in the argument of $F$, as $z$
traverses $\Gamma^+$.  Let $\Delta_\ell$ equal $\Delta \arg F(z)$ as
$z$ traverses $\Gamma_\ell$.  For $R$ sufficiently large the term
$a_N z^{N-k-1}$ in (\ref{eq:6}) dominates so that $\Delta_1 = (N-k-1)
\theta_0 +O (\frac{1}{R})$.  For $r$ sufficiently small and positive
the term $a_0 z^{-k-1}$ in (\ref{eq:6}) dominates so that $\Delta_5 =
-(-k-1) \theta_0 + O(r) = (k+1) \theta_0 + O(r)$, noting the
clockwise transversal of $\Gamma_5$.  Hence, \begin{equation}
\Delta_1 + \Delta_5 = N\theta_0 + O \left(\frac{1}{R}\right) +
O(r)\label{eq:7}\end{equation} for $1/R$ and $r$ sufficiently small.
Let $\Gamma_6 = \{t e^{i\theta_0}: r \leq t \leq R\}$, noting that
$\Gamma_2 \cup \Gamma_3 \cup \Gamma_4$ approaches $\Gamma_6$ as $s
\rightarrow O$.  For notational convenience let $\hat{\Delta}_6$
denote $\Delta_2+\Delta_4 +\lim \Delta_3$ as $s\rightarrow0$.  We observe that as $t$ crosses $1$ the change in
argument of $-\pi$ accounts for the change in argument determined by
traversing $\Gamma_3$ for $s$ sufficiently small, i.e., $\Delta_3 =
-\pi + O(s)$.  

Since a bound for $\hat{\Delta}_6$ can be determined by the
maximum number of times the image of $\Gamma_6$ crosses the real axis
we consider $\alpha(f_0)$, the number of real positive zeros of $f_0$
defined by \begin{eqnarray} f_0 (t) & = & Im \left[t^{k+1} F (te^{i
\theta_{0}}) \right]\label{eq:8}\\[2ex] & = & -\sum^k_{\ell=0}
a_\ell \left[\sin (k+1-\ell) \theta_0\right] t^\ell + \sum^N_{\ell =
k+2} a_\ell \left[\sin (\ell -k-1)\theta_0\right]
t^\ell\nonumber.\end{eqnarray} Descartes' Rule of Signs says that
$\alpha (f_0)$ is bounded above by $\gamma (f_0)$, the number of sign
changes in the coefficients of $f_0$.  Since, $a_\ell \geq 0$ for
each $\ell,\gamma(f_0)$ is determined by the number of sign changes
in $\sin \ell\theta_0$ as $\ell$ goes from $-k-1$ to $N-k-1$, i.e.,
over a range of $N$.  It is clear that there exists an $m$ such that
\begin{equation} \frac{m\pi}{N} < \theta_0 \leq \frac{(m+1)\pi}{N}.\label{eq:9}\end{equation} Since the number of sign changes as
$\ell\theta_0$ varies over an $N \theta_0$ range is determined by
$N\theta_0 /\pi$, it follows that $\gamma (f_0) \leq m+1$.  In
counting $\alpha(f_0)$ and  $\gamma(f_0)$ if $a_0 \sin(k+1)\theta_0 = 0$, then we factor out the leading power of $t^\ell$ in $f_0$.  So,
we are only considering positive zeros of $f_0$.

In the case when $\alpha(f_0)$ is strictly less than $m$, say $m_1$
with $0\leq m_1 <m$, it follows that $\hat{\Delta}_6 <m_1\pi +\pi
\leq m\pi$.  Thus, by using (7), (8) and (9) we obtain
$$\Delta\arg F(z)\geq 2\sum^5_{\ell=1}>2(N\theta_0-(m_1+1)\pi)\geq
0.$$  This would imply that $\Gamma$ encloses a zero of $F$,
contradicting the construction of $\Gamma$.

Now we consider the remaining cases, i.e., when $\gamma(f_0)=m$ or
$m+1$ and $\alpha(f_0) \geq m$.  If $\gamma(f_0)-\alpha (f_0)=1$,
then one of  the positive zeros of $f_0$ occurs with even
multiplicity.  This can be easily seen by observing that the sign of
$f_0(t)$ for large $t$ can be determined by  multiplying the sign of
the constant term of $f_0$ by $(-1)^{\gamma(f_0)}$ and, alternatively,
by multiplying the sign of the constant term of $f_0$ by
$(-1)^{\alpha(f_0)}$ if all of the positive zeros of $f_0$ have odd
multiplicity.  But, a positive zero of $f_0$ with even multiplicity
does not correspond to a proper crossing of the real axis by
$f(te^{i\theta_0})$ and, hence, $\hat{\Delta}_6\leq m\pi$ in this
case, which is impossible.

The only remaining case is when $\alpha(f_0) = \gamma(f_0)=m$ or
$m+1$.   We shall need an involved technical lemma, which we state and prove asLemmas 3.3 and 3.4 in the next section. To complete the proof we
state Lemma 3.3 as Lemma 2.1 in the following form:

\begin{lemma}  Let $g$ and $h$ be two real polynomials with $h(z) =
a_{\ell} z^{\ell} + \cdots + a_{n} z^{n} $ where $a_{\ell} > 0$ and
the degree of $g$ is less than $\ell$ with the last coefficient of
$g$ being negative.  If $p$ is the polynomial defined by $p = g + h$
and the number of sign charges in the coefficients of $p$ equals
the number of real positive zeros of $p$ then $h (z_{j})$ is positive
at each of the zeros, $z_{j},$ of $p$.

\end{lemma}

 We observe that the function $f_0$ defined in (\ref{eq:8}) satisfies
the hypothesis of Lemma 2.1 with $f_0 = g + h$ where $g(t) = -
\displaystyle{\sum^{k}_{l=1}} a_{\ell} \left[ \mbox{ sin } (k + 1 -
\ell ) \theta_{0} \right] t^{\ell}$ and $h(t) =
\displaystyle{\sum^{N}_{l = k+2}} a_{\ell} \left[\mbox{ sin } ( l -
 k- 1 ) \theta_{0} \right] t^{\ell}.$  Since $\alpha (f_0) = \gamma
(f_0)=m$ or $m+1$, it  follows, from a partitioning argument, that
the last nonzero term of $g(t)$ is negative and that the first term
of $h(t)$ is positive, for otherwise we would have $\gamma(f_0)\leq
m-1$, which would reduce to the previous case.  Now, since
$f_0(1)=0$, Lemma 2.1 gives that $h(1) > 0$, which
implies that $g(1) < 0$.  But, using (\ref{eq:3}) it follows that
$g(1) = -b_k \sin \theta_0$, which contradicts the assumption in
(\ref{eq:5}).  Our main theorem follows.

\section{Lemmas on Zeros, Critical Points, and Sign Changes} Let
$f(x)$ denote a polynomial with real coefficients.  Write $$f(x) =
\pm \left(p_0 (x) - p_1 (x) + p_2 (x) - p_3 (x) + \cdots \right)$$
where \begin{eqnarray*} p_i (x) & = & c_{i0} x^{n_{i-1}} + c_{i1}
x^{n_{i-1}+1} + \cdots + c_{i,m_i} x^{n_{i-1} +m_i}\\[2ex] n_{-1} & =
& 0,\ n_i = n_{i-1} + m_i +1,\\[2ex] c_{ij} & \geq & 0\ , i \geq 0 ,
j \geq 0.\end{eqnarray*}\vspace{12pt}

\noindent {\bf Notation}  {\em $\alpha (f)$ denotes the number of strictly
positive real zeros, $\gamma (f)$ denotes the number of sign changes,
$\tilde{\beta}(f)$ denotes the number of strictly positive critical
points,
$$\beta(f) = \left\{\begin{array}{cll}
0 & {\rm if} & f \equiv c_{00}\\[2ex]
\tilde{\beta} (f) + 1 & {\rm if} & p_0 (x) \equiv c_{00}\\[2ex]
\tilde{\beta} (f) & {\em otherwise}. &\end{array}\right.$$
Let $0 < x_1 (f) < x_2 (f) < x_3 (f) < \cdots $ denote the
distinct positive real zeros of $f$.  Let $y_i (f)$ denote the
critical point between $x_i(f)$ and $x_{i+1} (f)$ (in the case where
$\alpha(f) = \gamma(f)$).}

\begin{lemma}  \begin{itemize}
\item[(i)] $\alpha(f) \leq \beta (f) \leq \gamma (f)$.
\item[(ii)] If $\alpha (f) = \gamma(f)$, then $\alpha(f^{(k)}) =
\gamma(f^{(k)})$ for all $k$ such that $f^{(k)} \not\equiv 0$.
\end{itemize}
\end{lemma}

\noindent {\bf Proof.} 

\begin{itemize} 

\item[(i)] This is an immediate consequence of Rolle's Theorem and
Descartes' Rule of Signs.

\item[(ii)] It needs only to be shown that $\alpha (f^{'} ) = \gamma
(f^{'})$

\end{itemize}

Case 1: $p_0(x) \neq c_{00}$.  Then $\alpha (f^\prime) =
\tilde{\beta}(f) = \beta(f) = \gamma (f) = \gamma (f^{'})$. 

Case 2: $p_0 (x) = c_{00}$.  Then, $\alpha (f^{'}) = \tilde{\beta}
(f) = \beta(f) - 1 = \gamma (f) - 1 = \gamma (f^{'})$ $\Box$


\begin{remark}  The remaining lemmas will each contain
the hypothesis that $\alpha (f) = \gamma (f)$.  It follows from Lemma
3.1(i) that $\alpha (f) = \beta (f)$.  Therefore, the positive real
zeros and critical points of $f$ must strictly interlace. We
emphasize again that all coefficients of each of the polynomials
$p_{i} (x)$ appearing below are nonnegative. It will be useful to
the reader at this stage to note that the essentially different
cases are,

\begin{enumerate}

\item[(i)] $p_{0} \not\equiv c_{00} > 0$

\item[(ii)] $p_{0} \equiv c_{00} > 0 $

\item[(iii)] $c_{00} = 0$

\end{enumerate} It also may be beneficial to draw rough sketches for
each of the cases.  \end{remark}

\begin{lemma}  Let $f(x) = p_0 (x) - p_1 (x) + h(x)$, where $h(x) =
p_2(x) - p_3 (x) + \cdots$.  Suppose that $\alpha (f) = \gamma (f)$,
then $h(x_i (f)) > 0$ for all $i$.\end{lemma}

\noindent {\bf Proof.}  Let $p_1 (x) = c_n x^n + c_{n+1} x^{n+1} +
\cdots + c_r x^r$, where $c_r > 0$ and $n=0$ is allowed.  (If $n=0$,
then $p_0 (x) \equiv 0$.)

Since the lemma is trivial if $\alpha (h) = 0$, let us assume that
$\alpha (h) > 0$.  Observe that $x_1 (h^{(j)})$ is defined for $0
\leq j \leq r+1$, since $h^{(j)} (0) =0$ for $0 \leq j \leq r$.  Now
$f^{(r)} (x) = -r! c_r + h^{(r)}$ has a critical point at $x_1
(h^{(r+1)})$.  Since $\alpha (f^{(r)}) = \gamma (f^{(r)})$ (Lemma
3.1(ii)), it follows that $x_1 (f^{(r)}) < x_1 (h^{(r+1)})) < x_1
(h^{(r)})$, and $f^{(r)} (x_1 (h^{(r+1)})) > 0$.  Furthermore, since
$f^{(r)} < h^{(r)}$ it follows that $x_1 (f^{(r)}) < x_1 (h^{(r+1)}) <
x_2 (f^{(r)}) < x_1 (h^{(r)})$.  Now $f^{(r-1)}$ has critical points
at $x_1(f^{(r)})$ and $x_2(f^{(r)})$ and since $\alpha (f^{(r-1)}) =
\gamma(f^{(r-1)})$ it follows that $x_1 (f^{(r)}) < x_1 (f^{(r-1)}) <
x_2(f^{(r)})$.  Therefore, $f^{(r-1)} (x_2(f^{(r)})) > 0$, and since
$f^{(r-1)} < h^{(r-1)},$ it now follows that $x_2(f^{(r)}) <
x_2(f^{(r-1)}) < x_1(h^{(r-1)}).$ Thus, $x_{1} ( f^{(r-1)} ) < x_{2}
( f^{(r)} ) < x_{2} ( f^{(r-1)} ) < x_{1} ( h^{(r-1)}).$  An easy
induction then shows that $$x_1(f^{(j)}) < x_2(f^{(j)}) <
x_1(h^{(j)}), n\leq j \leq r-1.$$ 

Now let $p_0 (x) = d_0 + d_1 x + \cdots + d_{n-1} x^{n-1}$. Since it
suffices to prove the lemma in the generic case, we assume that
$d_{i} > 0 ; \ 0 \leq i \leq n-1.$ Since $f^{(n-1)}$ has critical
points at $x_1 (f^{(n)})$ and $x_2 (f^{(n)})$ and $\alpha (f^{(n-1)})
= \gamma (f^{(n-1)})$, it follows that $f^{(n-1)} (x_1 (f^{(n)})) <
0$ and $f^{(n-1)} (x_2 (f^{(n)})) > 0$.  Since $f^{(n)} (x) < h^{(n)}
(x)$ for $x > 0$, it now follows that for all $x \geq x_1 (f^{(n)})$,
\begin{eqnarray*} f^{(n-1)} (x) &=& f^{(n-1)} \left(x_1
\left(f^{(n)}\right)\right) + \int^x_{x_1 \left(f^{(n)}\right)}
f^{(n)} (t) dt\\[2ex] &<& \int^x_{x_1 \left(f^{(n)}\right)} h^{(n)}
(t) dt < h^{(n-1)} (x).\end{eqnarray*}

By a previous argument, we have that $x_1 (f^{(n)}) < x_2 (f^{(n)}) <
x_1 (h^{(n)}) < x_1 (h^{(n-1)})$.  Then, since $f^{(n-1)}  (x) <
h^{(n-1)} (x)$ for $x \geq x_1 (f^{(n)})$ and since $f^{(n-1)} (x_2
(f^{(n)})) > 0$, we have that $x_3 (f^{(n-1)}) < x_1 (h^{(n-1)})$.  Therefore,
$$\begin{array}{c}
\displaystyle{x_1 \left(f^{(n-1)}\right) < x_1
\left(f^{(n)}\right) < x_2 \left(f^{(n-1)}\right) < x_2
\left(f^{(n)}\right)}\\[2ex]
\displaystyle{< x_3 \left(f^{(n-1)}\right) < x_1
\left(h^{(n-1)}\right).}\end{array}$$\vspace{12pt}

Now suppose that for some $j$ with $0 < j \leq n-1$, $$x_1 (f^{(j)})
< x_2 (f^{(j)}) < x_3 (f^{(j)}) < x_1 (h^{(j)})$$ and $f^{(j)} (x) <
h^{(j)} (x) \mbox{ for all } x > y_1 (f^{(j)})$.

Then, exactly the same reasoning as above applies and it follows that
$$x_1 \left(f^{(j-1)}\right) < x_2 \left(f^{(j-1)}\right) < x_3
\left(f^{(j-1)}\right) < x_1 \left(h^{(j-1)}\right) \ \ \ \ \ \ \ \
(*^{j}) $$ and $f^{(j-1)} (x) < h^{(j-1)} (x)$ for all $x > y_1
(f^{(j-1)})$.  Therefore, $(\ast^j)$ is true for $0 \leq j \leq n-1$.


In particular, $(\ast^0)$ shows that $x_1 (f) < x_1 (h)$.  Thus,
$-p_0 (x_1 (f)) + p_1 (x_1(f)) = h(x_1(f)) > 0$.  By Lemma 3.1 $p_1 - 
p_0$ has only one zero. Therefore, $(p_1 - p_0)(x)$ is monotone
increasing for all $x > x_1 (p_1 - p_0)$.  Since $(p_1 - p_0) (x_1
(f)) > 0,$ we have $x_1 (f) > x_1 (p_1 - p_0)$ and it follows that $h(x_j
(f)) > 0$ for all $j$. $\Box$

\begin{lemma} Let $f = (p_0 - p_1) + (p_2 - p_3) + \cdots + (p_{2k+2}
- p_{2k+3}) + (p_{2k+4} - p_{2k+5}) + (p_{2k+6} - \cdots )$, where
$g_0 = p_0 - p_1,\ g_1 = p_2 - p_3, \cdots, g_{k+1} = p_{2k+2} -
p_{2k+3},\ h = p_{2k+4} - p_{2k+5} + p_{2k+6} - \cdots$.  Suppose
that $\alpha (f) = \gamma (f)$, then $h(x_i (f)) > 0$ for all
$i$.\end{lemma}

\noindent {\bf Proof.}  We may assume that $\gamma (h) \neq 0$, for
otherwise the result is trivial.  Hence, $\gamma (h^{(j)}) \neq 0$
for $0 \leq j \leq \deg (p_{2k+3})$.  Let $n_j = \deg (g_j)$
and $m_j = n_j +1$ for $0 \leq j \leq k+1$.  Let $m_{-1} =0$.

We now regard $k$ as fixed and proceed by induction on $\ell$, where
$-1 \leq \ell \leq k$.  From the previous lemma it follows that 
$$x_1 (f^{(m_k)}) < x_2 (f^{(m_k)}) < x_3 (f^{(m_k)}) < x_1 (h^{(m_k)})$$
and $(h^{(m_k)} - f^{(m_k)})(x)$ is positive and
monotone increasing for $x \geq y_1
(f^{(m_k)})$

We now make the following induction hypothesis: For some $\ell$ with
$0 \leq \ell \leq k$, $$x_1 (f^{(m_\ell)}) < \cdots < x_{2(k-\ell)
+3} (f^{(m_\ell)}) < x_1 (h^{(m_\ell)})$$ and $(h^{(m_\ell)} -
f^{(m_\ell)}) (x)$ is positive and monotone increasing for $x \geq
y_{2(k-\ell)+1} (f^{(m_\ell)})$.

Let $p_{2\ell+1} (x) = a_0 x^s + a_1 x^{s+1} + \cdots + a_r x^r$,
where $r = n_\ell = m_\ell -1$.  Now $f^{(r)} (x) = {\int^x_0}
f^{(m_\ell)} (t) dt - c$, where $c$ is a nonnegative constant.
Observe that $f^{(r)} (x)$ is positive for \begin{equation}
x_{2(k-\ell)+3} (f^{(r)}) < x < x_{2(k-\ell)+4}
(f^{(r)}).\label{eq:10}\end{equation} Therefore, for $x$ satisfying (10)
\begin{eqnarray*} 0 \leq f^{(r)} (x) & = & \int^x_{x_{2(k-\ell)+3}}
f^{(m_\ell)} (t) dt\\[2ex] & \leq & \int^x_{x_{2(k-\ell) +3}}
h^{(m_\ell)} (t) dt\end{eqnarray*} (by the induction hypothesis since
$x_{2(k-\ell)+3} (f^{(r)}) > y_{2(k-\ell)+1}(f^{(m_\ell)}))$ $$<
h^{(r)}(x_{2(k-\ell)+3}(f^{(r)})) + \int^x_{x_{2(k-\ell)+3}}
h^{(m_\ell)} (t) dt.$$

Therefore, 
$x_{2(k-\ell)+4} (f^{(r)}) < x_1
(h^{(r)})$ 
and since 
$\alpha (f^{(r)}) = \gamma (f^{(r)}),$
it follows that
\begin{equation}\begin{array}{llll}x_1(f^{(r)}) &<
x_1(f^{(m_\ell)}) &< x_2(f^{(r)}) &< \cdots \\[2ex]
&< x_{2(k-\ell)+3}(f^{(m_\ell)}) &< x_{2(k-\ell)+4} (f^{(r)}) &< x_1 (h^{(r)})
\end{array}
\label{eq:11}\end{equation}

and $(h^{(r)} - f^{(r)}) {(x)}$ is positive and monotone increasing
for $x > y_{2(k-\ell)+2} (f^{(r)})$.


A repetition of the above argument shows that $(11)$ remains true
when $r$ is replaced by $j$ and $m_\ell$ by $(j+1)$ for $s \leq j
\leq r$ and also shows that $x_{2(k-\ell)+5} (f^{(m_{\ell-1})}) < x_1
(h^{(m_\ell-1)})$.  This verifies the induction hypothesis for $(\ell
-1)$.

In particular, when $\ell = -1$ we have $x_1 (f) < x_2 (f) < \cdots <
x_{2k+5} (f) < x_1 (h)$ and $(h-f) (x) > 0$ for $x \geq y_{2k+3}
(f)$.  Since $h(x) > 0$ for $0 < x < x_1 (h)$ and since
$\sum^{k+1}_{j=0} g_j (x) <0$ for $x \geq y_{2k+3} (f)$ and $y_{2k+3}
(f) < x_1 (h)$, we finally have $h(x_i (f)) > 0$ for all $i$.
$\Box$\vspace{12pt}

The same argument establishes

\begin{lemma} Let $f = -p_0 + (p_1 - p_2) + \cdots + (p_{2k-1} -
p_{2k}) + (p_{2k+1} - p_{2k+2}) + \cdots$  Let $h = p_{2k+1} -
p_{2k+2} + \cdots$.  Suppose that $\alpha (f) = \gamma (f)$, then
$h(x_i (f)) > 0$ for all $i$.
\end{lemma}

\vspace{12pt}

\noindent{\bf \Large Acknowledgements}

\vspace{12pt}

  We wish to thank the mathematics department at the University of California at San Diego and the Math Sciences Research Institute at Berkeley for supporting the first and last authors respectively as visiting scholars during some of the preparation time of this paper.  We also wish to thank the number theorist
R. Evans of the University of San Diego for many helpful discussions.
\begin{thebibliography}{99}

\bibitem{1} Beauzamy, B. {\em Jensen's inequality for polynomials
with concentration at low degrees.}  Numer. Math. No. 49 (1986) pp.
221-225.

\bibitem{2} Beauzamy, B. and Enflo, P.  {\em Estimations de produits
de polyn\^{o}mes.} J. Number Theory. No. 21 (1985) pp. 390-412.

\bibitem{3} Evans, R. and Greene, J. {\em Polynomials with positive
coefficients whose zeros have modulus one.} preprint.

\bibitem{6} Evans, R. and Montgomery, P.  {\em Elementary Problem
Proposal for American Mathematical Monthly}.  Amer. Math. Mon. to appear.

\bibitem{4} Polya, G. and Szego, G. {\em Problems and theorems in
analysis II.} Springer-Verlag, 1976.

\bibitem{5} Rigler, A.K., Trimble, S.Y., and Varga, R.S. {\em Sharp
lower bounds for a generalized Jensen inequality.}  preprint.


\end{thebibliography}

\end{document}



